\chapter{2017年山东卷（理）}

\section{单选题}

\begin{problem}
设函数 \(y=\sqrt{4-x^2}\) 的定义域为 \(A\)，函数 \(y=\ln(1-x)\) 的定义域为 \(B\)，则 \(A\cap B=\)（\quad）.

\choices
    {\((1,2)\)}
    {\((1,2]\)}
    {\(\left(-2,1\right)\)}
    {\(\left[-2,1\right)\)}
\end{problem}

\begin{problem}
已知 \(a\in\R\)，\(\i\) 是虚数单位，若 \(z=a+\sqrt3\,\i\)，\(z\bar z=4\)，则 \(a=\)（\quad）.

\choices
    {\(1\) 或 \(-1\)}
    {\(\sqrt7\) 或 \(-\sqrt7\)}
    {\(-\sqrt3\)}
    {\(\sqrt3\)}
\end{problem}

\begin{problem}
已知命题 \(p:\forall x>0\)，\(\ln(x+1)>0\)；命题 \(q\)：若 \(a>b\)，则 \(a^2>b^2\)，下列命题为真命题的是（\quad）.

\choices
    {\(p\land q\)}
    {\(p\land \lnot q\)}
    {\(\lnot p\land q\)}
    {\(\lnot p\land \lnot q\)}
\end{problem}

\begin{problem}
已知 \(x\)，\(y\) 满足约束条件 \(x-y+3\le0\)，\(3x+y+5\le0\)，\(x+3\ge0\)，则 \(z=x+2y\) 的最大值是（\quad）.

\choices
    {\(0\)}
    {\(2\)}
    {\(5\)}
    {\(6\)}
\end{problem}

\begin{problem}
为了研究某班学生的脚长 \(x\)（单位：厘米）和身高 \(y\)（单位：厘米）的关系，从该班随机抽取 10 名学生，根据测量数据的散点图可以看出 \(y\) 与 \(x\) 之间有线性相关关系，设其回归直线方程为 \(\hat y=bx+a\)，已知 \(\sum_{i=1}^{10}x_i=225\)，\(\sum_{i=1}^{10}y_i=1600\)，\(b=4\)，该班某学生的脚长为 24，据此估计其身高为（\quad）.

\choices
    {\(160\)}
    {\(163\)}
    {\(166\)}
    {\(170\)}
\end{problem}

\begin{problem}
执行两次如图所示的程序框图，若第一次输入的 \(x\) 值为 7，第二次输入的 \(x\) 值为 9，则第一次，第二次输出的 \(a\) 值分别为（\quad）.

\bitmapfigure[float=inline,max width=0.42\linewidth]{img/2017/shandong_science/q6_fig1.png}

\choices
    {\(0,0\)}
    {\(1,1\)}
    {\(0,1\)}
    {\(1,0\)}
\end{problem}

\begin{problem}
若 \(a>b>0\)，且 \(ab=1\)，则下列不等式成立的是（\quad）.

\choices
    {\(a+\dfrac1b<\dfrac{b}{2a}<\log_2(a+b)\)}
    {\(\dfrac{b}{2a}<\log_2(a+b)<a+\dfrac1b\)}
    {\(a+\dfrac1b<\log_2(a+b)<\dfrac{b}{2a}\)}
    {\(\log_2(a+b)<a+\dfrac1b<\dfrac{b}{2a}\)}
\end{problem}

\begin{problem}
从分别标有 \(1\)，\(2\)，\(\ldots\)，\(9\) 的 9 张卡片中不放回地随机抽取 2 次，每次抽取 1 张，则抽到在 2 张卡片上的数奇偶性不同的概率是（\quad）.

\choices
    {\(\dfrac{5}{18}\)}
    {\(\dfrac{4}{9}\)}
    {\(\dfrac{5}{9}\)}
    {\(\dfrac79\)}
\end{problem}

\begin{problem}
在 \(\triangle ABC\) 中，角 \(A\)，\(B\)，\(C\) 的对边分别为 \(a\)，\(b\)，\(c\)，若 \(\triangle ABC\) 为锐角三角形，且满足 \(\sin B(1+2\cos C)=2\sin A\cos C+\cos A\sin C\)，则下列等式成立的是（\quad）.

\choices
    {\(a=2b\)}
    {\(b=2a\)}
    {\(A=2B\)}
    {\(B=2A\)}
\end{problem}

\begin{problem}
已知当 \(x\in[0,1]\) 时，函数 \(y=(mx-1)^2\) 的图象与 \(y=\sqrt{x}+m\) 的图象有且只有一个交点，则正实数 \(m\) 的取值范围是（\quad）.

\choices
    {\((0,1]\cup[2\sqrt3,+\infty)\)}
    {\((0,1]\cup[3,+\infty)\)}
    {\((0,\sqrt2)\cup[2\sqrt3,+\infty)\)}
    {\((0,\sqrt2]\cup[3,+\infty)\)}
\end{problem}

\section{填空题}

\begin{problem}
已知 \((1+3x)^n\) 的展开式中含有 \(x^2\) 的系数是 54，则 \(n=\fillinblank{}\).
\end{problem}

\begin{problem}
已知 \(\bs e_1\)，\(\bs e_2\) 是互相垂直的单位向量，若 \(\sqrt3\,\bs e_1-\bs e_2\) 与 \(\bs e_1+\lambda\bs e_2\) 的夹角为 \(60^\circ\)，则实数 \(\lambda\) 的值是\fillinblank{}.
\end{problem}

\begin{problem}
由一个长方体和两个 \(\dfrac14\) 圆柱构成的几何体的三视图如图，则该几何体的体积为\fillinblank{}.

\bitmapfigure[max width=0.48\linewidth]{img/2017/shandong_science/q13_fig1.png}
\end{problem}

\begin{problem}
在平面直角坐标系 \(xOy\) 中，双曲线 \(\dfrac{x^2}{a^2}-\dfrac{y^2}{b^2}=1\)（\(a>0\)，\(b>0\)）的右支与焦点为 \(F\) 的抛物线 \(x^2=2py\)（\(p>0\)）交于 \(A\)，\(B\) 两点，若 \(\lvert AF\rvert+\lvert BF\rvert=4\lvert OF\rvert\)，则该双曲线的渐近线方程为\fillinblank{}.
\end{problem}

\begin{problem}
若函数 \(\e^x f(x)\)（\(\e\approx2.71828\ldots\) 是自然对数的底数）在 \(f(x)\) 的定义域上单调递增，则称函数 \(f(x)\) 具有 \(M\) 性质. 下列函数中所有具有 \(M\) 性质的函数的序号为\fillinblank{}.
\begin{circlelist}
\item \(f(x)=2^{-x}\)
\item \(f(x)=3^{-x}\)
\item \(f(x)=x^3\)
\item \(f(x)=x^2+2\)
\end{circlelist}
\end{problem}

\section{解答题}

\begin{problem}
设函数 \(f(x)=\sin\!\left(\omega x-\dfrac{\pi}{6}\right)+\sin\!\left(\omega x-\dfrac{\pi}{2}\right)\)，其中 \(0<\omega<3\)，已知 \(f\!\left(\dfrac{\pi}{6}\right)=0\).

\begin{enumerate}
\item 求 \(\omega\)；
\item 将函数 \(y=f(x)\) 的图象上各点的横坐标伸长为原来的 2 倍（纵坐标不变），再将得到的图象向左平移 \(\dfrac{\pi}{4}\) 个单位，得到函数 \(y=g(x)\) 的图象，求 \(g(x)\) 在 \(\left[-\dfrac{\pi}{4},\dfrac{3\pi}{4}\right]\)上的最小值.
\end{enumerate}
\end{problem}

\begin{problem}
如图，几何体是圆柱的一部分，它是由矩形 \(ABCD\)（及其内部）以 \(AB\) 边所在直线为旋转轴旋转 \(120^\circ\) 得到的，\(G\) 是 \(\overset{\frown}{DF}\) 的中点.

\begin{enumerate}
\item 设 \(P\) 是 \(EC\) 上的一点，且 \(AP\perp BE\)，求 \(\angle CBP\) 的大小；
\item 当 \(AB=3\)，\(AD=2\) 时，求二面角 \(E-AG-C\) 的大小.
\end{enumerate}

\bitmapfigure[max width=0.55\linewidth]{img/2017/shandong_science/q17_fig1.png}
\end{problem}

\begin{problem}
在心理学研究中，常采用对比试验的方法评价不同心理暗示对人的影响，具体方法如下：将参加试验的志愿者随机分成两组，一组接受甲种心理暗示，另一组接受乙种心理暗示，通过对比这两组志愿者接受心理暗示后的结果来评价两种心理暗示的作用，现有 6 名男志愿者 \(A_1\)，\(A_2\)，\(A_3\)，\(A_4\)，\(A_5\)，\(A_6\) 和 4 名女志愿者 \(B_1\)，\(B_2\)，\(B_3\)，\(B_4\)，从中随机抽取 5 人接受甲种心理暗示，另 5 人接受乙种心理暗示.

\begin{enumerate}
\item 求接受甲种心理暗示的志愿者中包含 \(A_1\) 但不包含 \(B_1\) 的概率；
\item 用 \(X\) 表示接受乙种心理暗示的女志愿者人数，求 \(X\) 的分布列与数学期望 \(E(X)\).
\end{enumerate}
\end{problem}

\begin{problem}
已知 \(\{x_n\}\) 是各项均为正数的等比数列，且 \(x_1+x_2=3\)，\(x_3-x_2=2\).

\begin{enumerate}
\item 求数列 \(\{x_n\}\) 的通项公式；
\item 如图，在平面直角坐标系 \(xOy\) 中，依次连接点 \(P_1(x_1,1)\)，\(P_2(x_2,2)\)，\(\ldots\)，\(P_{n+1}(x_{n+1},n+1)\) 得到折线 \(P_1P_2\cdots P_{n+1}\)，求由该折线与直线 \(y=0\)，\(x=x_1\)，\(x=x_{n+1}\) 所围成的区域的面积 \(T_n\).
\end{enumerate}

\bitmapfigure[max width=0.48\linewidth]{img/2017/shandong_science/q19_fig1.png}
\end{problem}

\begin{problem}
已知函数 \(f(x)=x^2+2\cos x\)，\(g(x)=\e^x\left( \cos x-\sin x+2x-2 \right)\)，其中 \(\e\approx2.71828\ldots\) 是自然对数的底数.

\begin{enumerate}
\item 求曲线 \(y=f(x)\) 在点 \((\pi,f(\pi))\) 处的切线方程；
\item 令 \(h(x)=g(x)-af(x)\)（\(a\in\R\)），讨论 \(h(x)\) 的单调性并判断有无极值，有极值时求出极值.
\end{enumerate}
\end{problem}

\begin{problem}
在平面直角坐标系 \(xOy\) 中，椭圆 \(E:\dfrac{x^2}{a^2}+\dfrac{y^2}{b^2}=1\)（\(a>b>0\)）的离心率为 \(\dfrac{\sqrt2}{2}\)，焦距为 2.

\begin{enumerate}
\item 求椭圆 \(E\) 的方程；
\item 如图，该直线 \(l:y=k_1x-\dfrac{\sqrt3}{2}\) 交椭圆 E 于 A，B 两点，C 是椭圆 E 上的一点，直线 \(OC\) 的斜率为 \(k_2\)，且 \(k_1k_2=\dfrac{\sqrt2}{4}\)，\(M\) 是线段 \(OC\) 延长线上一点，且 \(\lvert MC\rvert:\lvert AB\rvert=2:3\)，\(\odot M\) 的半径为 \(\lvert MC\rvert\)，\(OS\)，\(OT\) 是 \(\odot M\) 的两条切线，切点分别为 \(S\)，\(T\)，求 \(\angle SOT\) 的最大值，并求取得最大值时直线 \(l\) 的斜率.
\end{enumerate}

\bitmapfigure[max width=0.55\linewidth]{img/2017/shandong_science/q21_fig1.png}
\end{problem}
